\documentclass[11pt,a4paper]{article}

\def\courseTitle{Industrial Security (Master)}
\def\sectionNumber{06}
\def\sectionTitle{Threat Intelligence for OT}
\def\sheetKind{Exercise Sheet}

% =====================================================================
% Shared preamble for exercise sheets and solution sheets.
%
% Define BEFORE \input{}-ing this file:
%   \def\courseTitle{...}
%   \def\sectionNumber{...}
%   \def\sectionTitle{...}
%   \def\sheetKind{Exercise Sheet}   or  Solution Sheet
%
% Optional:
%   \def\courseAuthor{...}
%
% The caller must issue \documentclass{...} (typically article) BEFORE
% loading this preamble.
% =====================================================================

\usepackage[utf8]{inputenc}
\usepackage[T1]{fontenc}
\usepackage[english]{babel}
\usepackage[a4paper, margin=2.2cm]{geometry}
\usepackage{graphicx}
\usepackage{listings}
\usepackage{xcolor}
\usepackage{tikz}
\usepackage{booktabs}
\usepackage{array}
\usepackage{enumitem}
\usepackage{titlesec}
\usepackage{fancyhdr}
\usepackage{hyperref}
\usepackage{amsmath}
\usepackage{amssymb}
\usepackage{tcolorbox}
\tcbuselibrary{skins, breakable}
\usepackage{needspace}

% Discourage solo lines split across pages.
\widowpenalty=10000
\clubpenalty=10000
\raggedbottom

% =====================================================================
% Shared TikZ styles for the Industrial Security lecture series.
% Loaded by every slide deck and exercise sheet that uses TikZ.
% =====================================================================

\usetikzlibrary{
  arrows.meta,
  shapes,
  shapes.geometric,
  shapes.symbols,
  shapes.multipart,
  positioning,
  fit,
  calc,
  backgrounds,
  decorations.pathmorphing,
  decorations.pathreplacing,
  decorations.text,
  patterns,
  matrix,
  shadows
}

% --- colour palette --------------------------------------------------
\definecolor{isOT}{HTML}{1F4E79}     % OT (deep blue)
\definecolor{isIT}{HTML}{6F2DA8}     % IT (purple)
\definecolor{isSafety}{HTML}{C00000} % safety red
\definecolor{isNet}{HTML}{2E7D32}    % network green
\definecolor{isAttack}{HTML}{B23A48} % attack red
\definecolor{isAccent}{HTML}{E08E0B} % amber
\definecolor{isMute}{HTML}{6B6B6B}   % grey
\definecolor{isLight}{HTML}{EDF2F8}  % light background
\definecolor{isPLC}{HTML}{2C5282}
\definecolor{isHMI}{HTML}{2F855A}
\definecolor{isSCADA}{HTML}{6B46C1}

% --- generic node styles --------------------------------------------
\tikzset{
  isnode/.style={
    draw, rounded corners=2pt, minimum height=8mm, minimum width=18mm,
    font=\small, align=center, thick
  },
  isbox/.style={isnode, fill=isLight},
  plc/.style={isnode, fill=isPLC!15, draw=isPLC, thick},
  hmi/.style={isnode, fill=isHMI!15, draw=isHMI, thick},
  scada/.style={isnode, fill=isSCADA!15, draw=isSCADA, thick},
  sensor/.style={isnode, fill=isNet!10, draw=isNet, minimum height=6mm, minimum width=14mm, font=\scriptsize},
  actuator/.style={isnode, fill=isAccent!15, draw=isAccent, minimum height=6mm, minimum width=14mm, font=\scriptsize},
  firewall/.style={isnode, fill=isAttack!15, draw=isAttack, thick},
  server/.style={isnode, fill=isMute!15, draw=isMute!80, thick},
  switch/.style={isnode, fill=isLight, draw=black!60, minimum height=6mm, minimum width=14mm, font=\scriptsize},
  workstation/.style={isnode, fill=white, draw=isOT, thick},
  attacker/.style={isnode, fill=isAttack!20, draw=isAttack, thick, font=\small\bfseries},
  inetnode/.style={draw, ellipse, minimum width=24mm, minimum height=10mm, font=\scriptsize, fill=isLight, thick},
  process/.style={isnode, fill=isOT!15, draw=isOT, thick, minimum width=24mm},
  zone/.style={
    draw, rounded corners=4pt, thick, dashed, inner sep=4mm
  },
  conduit/.style={-{Stealth[length=2.5mm]}, thick, isNet!80!black},
  attack/.style={-{Stealth[length=2.5mm]}, thick, isAttack, dashed},
  flow/.style={-{Stealth[length=2.5mm]}, thick, isOT!70!black},
  netline/.style={thick, black!60},
  axisline/.style={-{Stealth[length=2mm]}, thick, black!70},
  tag/.style={font=\scriptsize\itshape, fill=white, inner sep=1pt},
  level/.style={
    draw, fill=isLight, minimum width=0.85\linewidth, minimum height=8mm,
    font=\small, anchor=west
  },
  brace/.style={decorate, decoration={brace, amplitude=4pt, raise=2pt}},
  legendbox/.style={draw, rounded corners=2pt, fill=isLight, inner sep=2mm, font=\scriptsize, align=left}
}

% --- handy macros ----------------------------------------------------
% \plcicon, \hmiicon etc as simple labelled boxes; useful inline.
\newcommand{\plcicon}[1][PLC]{\tikz[baseline=-0.6ex]{\node[plc, minimum width=10mm, minimum height=5mm, font=\scriptsize, inner sep=1pt]{#1};}}
\newcommand{\hmiicon}[1][HMI]{\tikz[baseline=-0.6ex]{\node[hmi, minimum width=10mm, minimum height=5mm, font=\scriptsize, inner sep=1pt]{#1};}}
\newcommand{\fwicon}[1][FW]{\tikz[baseline=-0.6ex]{\node[firewall, minimum width=10mm, minimum height=5mm, font=\scriptsize, inner sep=1pt]{#1};}}


\providecommand{\courseAuthor}{Industrial Security Lecture Series}
\providecommand{\courseInstitute}{Department of Computer Science}
\providecommand{\companyLogo}{logo}

% --- Corporate branding overrides ------------------------------------
% Picks up logo / author / brand colours from the repo-root branding.tex
% when present; falls back to the defaults above otherwise.
\IfFileExists{../../../branding.tex}{% =====================================================================
% Corporate / institutional branding -- single file to customise the
% look of the entire course (slides, notes, exercises, labs).
%
% HOW TO REBRAND IN UNDER A MINUTE:
%   1. Replace `branding/logo.pdf` with your own logo
%      (PDF preferred; PNG and JPG also work -- name it `logo.<ext>`).
%   2. (Optional) change the two colours below to your brand palette.
%   3. (Optional) override the author/institute lines below.
%   4. Re-run `make` at the repo root. That's it.
%
% Everything in this file has sensible defaults; you can delete a line
% to fall back to the built-in defaults, or comment it out.
%
% This file is loaded automatically by the slides, exercise, and lab
% preambles -- you never need to \input it manually.
% =====================================================================

% --- Logo --------------------------------------------------------------
% Path is resolved from each leaf-build directory; the default points at
% the shared `branding/` folder at the repo root. If you put your logo
% somewhere else, set the full or relative path here (without extension):
\renewcommand{\companyLogo}{../../../branding/logo}

% --- Author / institute (appears on the title page footer) ------------
\renewcommand{\courseAuthor}{}
\renewcommand{\courseInstitute}{}

% --- Brand colours ----------------------------------------------------
% slideAccent      = primary brand colour (title bands, accent rule)
% slideSecondary   = secondary brand colour (section badge, highlight)
% Both use HTML hex codes. Keep enough contrast against white text.
\definecolor{slideAccent}{HTML}{00205B}      % deep navy
\definecolor{slideSecondary}{HTML}{00A0DC}   % light blue accent
}{}

% --- listings -------------------------------------------------------
\definecolor{codebg}{rgb}{0.96,0.96,0.96}
\lstset{
  backgroundcolor=\color{codebg},
  basicstyle=\ttfamily\footnotesize,
  breaklines=true,
  frame=single,
  numbers=left,
  numberstyle=\tiny\color{black!50},
  showstringspaces=false,
  tabsize=2
}

% --- header / footer ------------------------------------------------
\pagestyle{fancy}
\fancyhf{}
\renewcommand{\headrulewidth}{0.4pt}
\renewcommand{\footrulewidth}{0pt}
\lhead{\small \courseTitle}
\rhead{\small Section \sectionNumber{} -- \sheetKind}
\cfoot{\thepage}

% --- task / solution boxes -----------------------------------------
% `before={...}` guards every task/solution box with \needspace so the
% box doesn't start with only one or two lines of breathing room at the
% bottom of a page. If less than ~9 lines remain, LaTeX bumps the box
% to the next page; otherwise it starts here and breaks normally if it
% runs long.
% Boxes are `breakable` so very long solutions don't blow the page,
% but `before={\needspace{14\baselineskip}}` pushes them to a fresh
% page whenever fewer than ~14 lines remain. That covers the vast
% majority of single-question splits in practice.
\newtcolorbox{taskbox}[1]{
  enhanced, breakable,
  colback=isLight, colframe=isOT,
  fonttitle=\bfseries\color{white},
  coltitle=white,
  colbacktitle=isOT,
  title={#1},
  arc=2pt, boxrule=0.6pt,
  left=2mm, right=2mm, top=2mm, bottom=2mm,
  before={\par\vspace{2mm}\needspace{14\baselineskip}\par},
  after={\par\vspace{2mm}}
}

\newtcolorbox{solutionbox}{
  enhanced, breakable,
  colback=isNet!5, colframe=isNet,
  fonttitle=\bfseries\color{white},
  coltitle=white,
  colbacktitle=isNet,
  title={Solution},
  arc=2pt, boxrule=0.6pt,
  left=2mm, right=2mm, top=2mm, bottom=2mm,
  before={\par\vspace{2mm}\needspace{14\baselineskip}\par},
  after={\par\vspace{2mm}}
}

\newtcolorbox{hintbox}{
  enhanced, breakable,
  colback=isAccent!10, colframe=isAccent,
  fonttitle=\bfseries\color{white}, coltitle=white, colbacktitle=isAccent,
  title={Hint},
  arc=2pt, boxrule=0.6pt
}

\newtcolorbox{warnbox}{
  enhanced, breakable,
  colback=isAttack!8, colframe=isAttack,
  fonttitle=\bfseries\color{white}, coltitle=white, colbacktitle=isAttack,
  title={Caution},
  arc=2pt, boxrule=0.6pt
}

% --- section formatting --------------------------------------------
\titleformat{\section}{\Large\bfseries\color{isOT}}{\thesection}{0.5em}{}
\titleformat{\subsection}{\large\bfseries\color{isOT!80!black}}{\thesubsection}{0.5em}{}

% --- title block macro ----------------------------------------------
\newcommand{\sheetheader}{%
  \begin{center}
    {\Large\bfseries \courseTitle}\\[0.3em]
    {\large \sheetKind{} -- Section \sectionNumber}\\[0.2em]
    {\large \sectionTitle}\\[0.2em]
    \rule{\linewidth}{0.4pt}
  \end{center}
  \vspace{0.5em}
}


\begin{document}
\sheetheader

\noindent
\textbf{Goal.} These exercises drill analytic tradecraft, not recall. You will build an
ATT\&CK-for-ICS profile for a named activity group, run a Diamond Model on a real incident,
grade intelligence by confidence, translate consequence into Priority Intelligence
Requirements (PIRs), turn a TTP into a deployable detection, dismantle an over-confident
attribution, and write a TLP-marked assessment. Assume the full Bachelor ICS background; no
fundamentals are re-taught here. Cite primary reporting (CISA, Dragos, Mandiant, Microsoft,
ESET, OASIS) and be explicit about what you do \emph{not} know. The lab (06) builds the
machine-readable artefacts (Navigator layer, STIX bundle); this sheet builds the reasoning
behind them -- align with it, do not duplicate it.

\begin{taskbox}{Exercise 6.1 -- ATT\&CK-for-ICS Technique Profile}
Choose \textbf{one} activity group and carry it through this whole sheet:
\textbf{ELECTRUM / Sandworm} (GRU Unit 74455), \textbf{CHERNOVITE} (PIPEDREAM /
INCONTROLLER author), or \textbf{VOLTZITE / Volt Typhoon}.
\begin{enumerate}\setlength{\itemsep}{0pt}
  \item Produce a technique profile of \textbf{at least six} distinct ATT\&CK \emph{for ICS}
        techniques (\texttt{T0xxx}) your group has demonstrably used. Present as a four-column
        table: \emph{Tactic, Technique~ID, Technique~name, Evidence sentence} (which report,
        which incident).
  \item Many real OT campaigns are \emph{dual-matrix}: enterprise (\texttt{T1xxx}) techniques
        for the IT intrusion, ICS (\texttt{T0xxx}) techniques for the OT effect. Name the one
        technique that marks the \textbf{IT$\to$OT pivot} for your group, and say why it is the
        pivot.
  \item State one technique you believe the group \emph{could} use but for which you have
        \textbf{no public evidence}, and mark it explicitly as a hypothesis, not a fact.
\end{enumerate}
\end{taskbox}

\begin{taskbox}{Exercise 6.2 -- Diamond Model of One Incident}
Pick \textbf{one} concrete incident attributed to your Ex.~6.1 group (e.g.\ Industroyer2 at a
Ukrainian transmission utility, April~2022; the 2015 Ukraine distribution blackout;
PIPEDREAM's pre-deployment discovery, 2022; VOLTZITE pre-positioning in US comms/electric,
2023--24). Fill all four vertices with \emph{specific} evidence, plus the meta-features:
\begin{enumerate}\setlength{\itemsep}{0pt}
  \item \textbf{Adversary} -- named group and, if known, sub-cluster or operator signature.
  \item \textbf{Capability} -- the specific tooling/malware and its ATT\&CK \texttt{Txxxx}
        IDs (cross-reference Ex.~6.1).
  \item \textbf{Infrastructure} -- C2 / staging / delivery, classed Type~1 (adversary-owned)
        vs.\ Type~2 (hijacked third party).
  \item \textbf{Victim} -- sector, asset class, and the \emph{process consequence} the
        adversary was reaching for.
  \item \textbf{Meta} -- timestamp, ICS Cyber Kill Chain stage (Stage~1 vs.\ Stage~2),
        result, and an explicit confidence rating with one sentence of justification.
\end{enumerate}
Cite a source for every vertex. State at least one vertex you had to leave as ``unknown''.
\end{taskbox}

\begin{taskbox}{Exercise 6.3 -- IOC vs.\ TTP vs.\ Strategic Intel, and Confidence}
You receive four artefacts about your group:
\begin{enumerate}\setlength{\itemsep}{0pt}
  \item A SHA-256 hash of a dropper, from one vendor blog, observed once.
  \item ``The group uses \texttt{netsh portproxy} to relay traffic through compromised SOHO
        routers'' -- corroborated by two independent vendors.
  \item A single forum post claiming the group is ``preparing to hit European LNG this
        winter''.
  \item ``This activity is consistent with state-directed pre-positioning for disruptive
        effect during a crisis'' -- a government assessment.
\end{enumerate}
For each: (a)~classify it as IOC / TTP / strategic intel; (b)~place it on the Pyramid of Pain;
(c)~assign a confidence using \emph{both} the Admiralty (NATO) source-reliability $\times$
information-credibility grade (A--F / 1--6) \emph{and} a Words of Estimative Probability band
(e.g.\ ``likely'', ``highly likely''); (d)~justify each grade in one sentence. Present as a
table.
\end{taskbox}

\begin{taskbox}{Exercise 6.4 -- Consequence-Driven PIRs for a Sector}
Pick a \emph{named} sub-sector your Ex.~6.1 group threatens (electric \emph{transmission},
water \emph{treatment}, or \emph{LNG} export). Write \textbf{five} ranked Priority
Intelligence Requirements. PIRs are \emph{questions}, not statements, and must be
consequence-driven: anchor each to a process consequence the operator fears (loss of view,
loss of control, equipment damage, safety trip), never a generic IT outcome. For each PIR
give: the question; the consequence it protects against; the collection source that would
answer it; and the trigger / threshold that escalates it. Then, in two sentences, justify your
\textbf{ranking} -- why is PIR-1 above PIR-5?
\end{taskbox}

\begin{taskbox}{Exercise 6.5 -- Map a TTP to a Detection}
Take one ICS-protocol or LotL technique from your Ex.~6.1 profile and turn it into a
\emph{deployable} detection. Either a \textbf{Sigma} rule (for an IT/log-based TTP) or a
\textbf{Zeek} script/notice (for an OT-protocol TTP). Your answer must include:
\begin{enumerate}\setlength{\itemsep}{0pt}
  \item the \texttt{Txxxx} ID the rule covers, embedded so the rule self-documents its mapping;
  \item the telemetry source it needs and the exact logic (field, condition, threshold/window);
  \item one realistic false-positive condition and a one-line tuning note that suppresses it
        without blinding the rule.
\end{enumerate}
Place any rule body in a \texttt{lstlisting}. State where this detection sits on the Pyramid of
Pain and what the adversary must do to evade it.
\end{taskbox}

\begin{taskbox}{Exercise 6.6 -- Critique an Over-Attribution Claim}
A junior analyst circulates this line in a draft report:

\begin{quote}\itshape
``We observed Modbus write bursts to a PLC on TCP/502 from an internal host, plus a
\texttt{netsh portproxy} command on a jump box. This is Volt Typhoon. Notify the board that a
PRC state actor is inside our OT network.''
\end{quote}

Identify \textbf{four} distinct analytic errors in this claim (e.g.\ base-rate / shared-TTP
problems, the Pyramid of Pain level used, missing the IT$\to$OT distinction, confidence
language, consumer fit). For each, state the error and the correction. Then rewrite the claim
into one defensible sentence that an analyst could actually stand behind, using calibrated
estimative language.
\end{taskbox}

\begin{taskbox}{Exercise 6.7 -- TLP-Marked One-Paragraph Assessment}
Synthesise your work into a single \textbf{TLP-marked} assessment paragraph (5--8 sentences)
that a SOC lead can read in under a minute. It must contain: a calibrated bottom-line-up-front
judgement (with estimative language and a confidence level); the one or two TTPs that matter
most this week; a named, prioritised recommended action tied to a consequence; and the correct
\textbf{TLP} marking for a product synthesised from public (TLP:CLEAR) reporting but intended
only for your internal cohort. Justify your TLP choice in one separate sentence after the
paragraph.
\end{taskbox}

\begin{taskbox}{Exercise 6.8 -- Essay: The Limits of IOC-Based Detection for LotL Actors}
Write a one-page (about 400--600~word) analytic essay on why indicator-based detection
structurally fails against a living-off-the-land actor such as VOLTZITE / Volt Typhoon, and
what replaces it. Your essay should: (i)~explain, using the Pyramid of Pain, \emph{why} the
indicators these actors leave are exactly the cheapest layer to change, and why their tooling
(signed Windows binaries, hijacked SOHO routers as Type-2 infrastructure) defeats hash- and
IP-based feeds; (ii)~name the detection classes that \emph{do} work (behavioural / parent--child
process analytics, identity / UEBA, OT-protocol baselining) and give one concrete example each;
(iii)~address the OT-specific twist -- why the IT-side LotL detections do not directly transfer
to a Purdue Level~1/2 network and what telemetry you would need there; (iv)~conclude with one
honest limitation of the behavioural approach (e.g.\ baseline cost, false-positive load,
detection latency). Cite at least one primary source (CISA \emph{AA24-038A}, Microsoft 2023, or
Dragos VOLTZITE reporting).
\end{taskbox}

\end{document}
